% Generated draft: requires mathematical and source review.
\documentclass[11pt]{article}
\usepackage[margin=25mm]{geometry}
\usepackage{fontspec}
\setmainfont{Latin Modern Roman}
\usepackage{amsmath,amssymb,mathtools,unicode-math}
\setmathfont{Latin Modern Math}
\usepackage{enumitem,needspace,xurl,hyperref}
\setlength{\parindent}{0pt}
\setlength{\parskip}{5pt}
\setlength{\emergencystretch}{3em}
\newcommand{\N}{\mathbb{N}}
\newcommand{\Z}{\mathbb{Z}}
\newcommand{\Q}{\mathbb{Q}}
\newcommand{\R}{\mathbb{R}}
\newcommand{\C}{\mathbb{C}}
\newcommand{\F}{\mathbb{F}}
\newcommand{\A}{\mathbb{A}}
\newcommand{\PP}{\mathbb{P}}
\newcommand{\E}{\mathbb{E}}
\newcommand{\eps}{\varepsilon}
\newcommand{\dd}{\,\mathrm{d}}
\begin{document}
\section*{UnsolvedMath ID: 701552}
\textbf{Source problem code:} catalogue entry
\subsection{Definitions and mathematical statement}
A group is finitely generated free if it has a finite free basis. A finitely generated free abelian group is a group isomorphic to \(\Z^n\) for some integer \(n\geq 0\). A cyclic group is a group generated by one element.

The supplied record asks the following two questions.
\begin{enumerate}
\item[(a)] Let \(I\) be an infinite index set and let \((F_i)_{i\in I}\) be a family of finitely generated free groups. If \(G\) is their free product with cyclic amalgamation, must
\[
\psi(G)=1
\]
hold?
\item[(b)] Let \(I\) be an infinite index set and let \((A_i)_{i\in I}\) be a family of finitely generated free abelian groups. If \(G\) is their free product with cyclic amalgamation, must
\[
\psi(G)=1
\]
hold?
\end{enumerate}

\textbf{Unresolved formulation qualifications.} The supplied record does not define \(\psi(G)\), so the displayed conclusions retain an undefined source symbol. In particular, interpreting \(1\) as a trivial subgroup would require knowing that \(\psi(G)\) denotes a subgroup. The record also does not specify the amalgamation construction for infinitely many factors: it does not identify a common cyclic subgroup and its embeddings, or describe an iterated amalgamation or an underlying graph of groups. It gives no explicit restrictions on the orders of the cyclic amalgamating groups, their possible triviality, the ranks of the factors, or the cardinality of the infinite family. Thus these questions preserve the supplied wording but do not determine fully specified mathematical assertions without additional source context.
\subsection{Short English statement}
The source asks whether \(\psi(G)=1\) for a free product of infinitely many finitely generated free groups with cyclic amalgamation, and asks the same question for finitely generated free abelian factors. The meaning of \(\psi\) and the intended infinite amalgamation construction are missing from the supplied record.
\subsection{Sources}
\begin{itemize}
\item UnsolvedMath local registry, permanent ID \(701552\), title \emph{Kourovka Notebook Problem 19.4}, exact source code \texttt{catalogue entry}; upstream ID \(701552\), upstream revision \texttt{d371ee917ac9647fbb06cd6fc529deec6f3dda82}. The supplied record was created and updated at \texttt{2026-10-16T14:00:00Z} and carries the status label \texttt{open}; this label is supplied registry metadata, not an independently verified mathematical assessment.
\item Supplied bibliographic attribution: \emph{The Kourovka Notebook (Unsolved Problems in Group Theory)}, 19th issue (2018), problem 19.4; also identified in the record with the 21st issue edition (2026), \texttt{arXiv:1401.0300v48}. Proposed by M. K. Azarian in 2018. The supplied collection description names E. I. Khukhro and V. D. Mazurov as editors of the Issues 1--21 collection (1965--2026).
\item The supplied bibliographic reference has not been independently consulted. No definition of \(\psi\), explanatory source excerpt, or literature-verification date is included in the supplied record.
\end{itemize}
\end{document}
